\chapter{The Mesh-LZ Architecture}
\label{ch:architecture}

The architecture of the Mesh-LZ codec is explicitly designed around block-level independence, deterministic spatial stenciling, and localized Lempel-Ziv parsing. By enforcing strictly independent block processing, the codec naturally lends itself to massive parallelization, satisfying the modern requirement for high-throughput multi-core execution on superscalar processors.

\section{System Pipeline Overview}

The Mesh-LZ encoding pipeline proceeds through a series of deterministic stages, transforming raw 2D pixel arrays into a densely packed, interleaved rANS bitstream.

\subsection{Encoding Flow}
\begin{enumerate}
    \item \textbf{Padding and Normalization:} The input image of width $W$ and height $H$ is padded using edge-clamping (repeating the edge pixels) to ensure the dimensions are strictly divisible by the target block size $b$ (typically $b=8$ or $b=16$). This eliminates complex bounds-checking during the inner loops of the compression engine.
    \item \textbf{Color Space Transformation (Optional):} If the encoder is configured for lossy or highly decorrelated compression, the raw 24-bit RGB data is mathematically converted to the reversible YCoCg-R color space.
    \item \textbf{Chroma Subsampling (Optional):} For web-optimized delivery modes, 4:2:0 chroma subsampling is subsequently applied. This is achieved by averaging $2 \times 2$ pixel clusters in the $Co$ and $Cg$ planes, reducing the spatial resolution of the color data by 75\% while leaving the Luma ($Y$) channel fully intact.
    \item \textbf{Palette Quantization (Optional):} If the \texttt{--palette} mode is active, the 24-bit TrueColor data is mapped to an optimal 256-color palette utilizing a Kohonen neural network vector quantization algorithm (NeuQuant). The image is subsequently processed as a single-channel 8-bit index map.
    \item \textbf{Dynamic Stenciling:} The image is subdivided into independent $b \times b$ blocks. For each block, the encoder calculates the discrete horizontal and vertical gradients to select an optimal 1D traversal stencil (Raster, Columnar, or Hilbert).
    \item \textbf{Lempel-Ziv Parsing:} The 1D pixel stream generated by the selected stencil is parsed using a greedy Lempel-Ziv dictionary matching algorithm, producing a sequence of literal pixels and \texttt{[offset, length]} backward references.
    \item \textbf{Residual Quantization (Optional):} Literal pixels undergo a predictive step where the difference (residual) between the current pixel and the previous decoded pixel is computed. For lossy compression, this residual is divided by a uniform scalar quantization step $q$. The residual is then zig-zag mapped to a positive integer.
    \item \textbf{Frequency Analysis:} The occurrences of all symbols (stencils, LZ commands, residuals, offsets, lengths) are counted across the entire image to construct global probability models.
    \item \textbf{rANS Entropy Coding:} The symbols are encoded in reverse order into multiple interleaved byte streams using the constructed frequency tables.
\end{enumerate}

\subsection{Decoding Flow}
The decoding pipeline mirrors these steps in reverse. Crucially, the decoder \textit{does not} need to perform the expensive gradient variance analysis or the string matching searches. It simply decodes the symbols via the interleaved rANS engine, recreates the 1D path according to the decoded stencil index, and reconstructs the spatial block by copying dictionary references. This fundamental asymmetry is what allows Mesh-LZ to decode in single-digit milliseconds.

\section{Dynamic Stencil Selection Heuristics}

The transformation of a 2D block into a 1D sequence is the most critical step in enabling the Lempel-Ziv algorithm to identify spatial redundancy. A static traversal path (e.g., a standard left-to-right Raster scan) fails spectacularly when images exhibit strong vertical features, such as the edge of a building or a tree trunk. If a vertical line is parsed using a horizontal raster scan, the pixels belonging to that line are separated by $b-1$ other pixels in the 1D stream, preventing the LZ parser from finding a contiguous match.

Mesh-LZ introduces a dynamic stenciling mechanism. For every block, the encoder calculates two gradient variance metrics, $V_H$ (Horizontal Variance) and $V_V$ (Vertical Variance):
\begin{equation}
    V_H = \sum_{y=0}^{b-1} \sum_{x=0}^{b-2} \big| P(x, y) - P(x+1, y) \big|
\end{equation}
\begin{equation}
    V_V = \sum_{y=0}^{b-2} \sum_{x=0}^{b-1} \big| P(x, y) - P(x, y+1) \big|
\end{equation}
where $P(x,y)$ represents the luminance intensity of the pixel at coordinate $(x,y)$.

A heuristic selection is then applied to choose the 1D path:
\begin{itemize}
    \item \textbf{Column Scan (Vertical bias):} If $V_H > 1.5 \times V_V$ and $V_H$ exceeds a basic noise threshold, the block exhibits strong horizontal variance. This implies the presence of vertical structural bands. The encoder selects the Column Scan stencil (traversing top-to-bottom, left-to-right).
    \item \textbf{Raster Scan (Horizontal bias):} If $V_V > 1.5 \times V_H$ and $V_V$ exceeds the noise threshold, the block exhibits strong vertical variance, implying horizontal structural bands. The encoder selects the Raster Scan stencil.
    \item \textbf{Hilbert Space-Filling Curve:} If the variances are relatively balanced ($|V_H \approx V_V|$), indicating complex localized texture or diagonal features, the encoder defaults to a Hilbert Curve. The Hilbert curve is mathematically proven to preserve 2D locality remarkably well in 1D space, ensuring that pixels physically close in 2D remain closely indexed in the 1D array.
\end{itemize}

By dynamically adapting the traversal path on a block-by-block basis, Mesh-LZ maximizes the probability that consecutive pixels in the 1D stream will be identical or highly correlated, drastically improving the LZ dictionary hit rate.

\section{Scalar Quantization and Zig-Zag Mapping}

For lossy encoding, Mesh-LZ employs a uniform scalar quantization scheme applied strictly to the spatial residuals. When the LZ parser encounters a pixel that cannot be matched in the localized dictionary, it encodes it as a literal. To exploit localized correlation, Mesh-LZ does not encode the absolute pixel value; instead, it encodes the residual (difference):
\begin{equation}
    R_{raw} = P_i - P_{i-1}
\end{equation}

In lossy modes, $R_{raw}$ is quantized by a globally defined step size $q$:
\begin{equation}
    R_{q} = \text{sign}(R_{raw}) \times \left\lfloor \frac{|R_{raw}| + q / 2}{q} \right\rfloor
\end{equation}

Because $R_q$ is a signed integer (representing both positive and negative differences), and the rANS entropy encoder strictly requires positive integer symbols starting from zero, Mesh-LZ utilizes a zig-zag mapping function to injectively map signed integers to unsigned integers:
\begin{equation}
    Z(x) = 
    \begin{cases} 
      2x & \text{if } x \ge 0 \\
      -2x - 1 & \text{if } x < 0 
   \end{cases}
\end{equation}

This mapping sequence operates as: $0 \rightarrow 0$, $-1 \rightarrow 1$, $1 \rightarrow 2$, $-2 \rightarrow 3$, $2 \rightarrow 4$, and so on. Because image residuals in natural photography follow a Laplacian or Cauchy distribution strongly geometrically centered around zero, this mapping ensures that smaller, vastly more probable residuals map to smaller unsigned integers. This concentrates the probability mass at the lower indices of the rANS frequency table, allowing the entropy coder to allocate the shortest fractional bit representations to the most common occurrences.
